%% kapa: exam=1 round=20 year=2009 period=1 form=A subject=economics subject_ko=경제원론 seq=24 number=64
%% source: Q-Net 감정평가사 기출문제 (cst00309) artlSeq=5024178 file=extracted/1008/5024178/2009년 제20회 감정평가사 1차시험 기출문제/1교시.hwp
%% conversion: kordoc 4.24.51 -> tools/gen1.py (auto); review=proofread-sonnet; flags=none
\begin{question}{64}
조세가 없는 경우 어떤 상품의 수요함수와 공급함수가 다음과 같다.
\begin{center}$Q_{ D } =500-0.5P$, $Q_{ S } =50+P$\end{center}
상품 한 단위당 30의 조세가 공급자에게 부과될 때 단위당 조세의 귀착은? (단, $Q_{ D }$ 는 수요량, $Q_{ S }$ 는 공급량, $P$ 는 가격이다.)

\noindent\hspace*{2em}\makebox[1.8cm][l]{\uline{수요자}}\uline{공급자}\par
\begin{choices}
  \item \makebox[1.8cm][l]{5}25
  \item \makebox[1.8cm][l]{10}20
  \item \makebox[1.8cm][l]{15}15
  \item \makebox[1.8cm][l]{20}10
  \item \makebox[1.8cm][l]{25}5
\end{choices}
\end{question}
\answer{4}
